\section{Evaluación y Benchmarking}
\label{sec:evaluacion}

La evaluación rigurosa de un modelo de lenguaje requiere ir más allá de la observación de la curva de pérdida durante el entrenamiento. Este capítulo documenta la evaluación cuantitativa del modelo híbrido tfm-slm mediante un benchmark exhaustivo de 10 métricas ejecutado sobre un conjunto de validación independiente, seguido de un análisis de los resultados obtenidos y su contexto respecto a modelos de referencia.

% ---------------------------------------------------------------
\subsection{Entorno y Configuración del Benchmark}

La evaluación se ejecuta sobre la instancia \texttt{g7e.2xlarge} equipada con una GPU \textbf{NVIDIA RTX PRO 6000 Server} (96 GB VRAM), el mismo hardware utilizado para el entrenamiento final. A lo largo del proyecto se emplearon tres generaciones de hardware de forma progresiva: la NVIDIA L4 (24 GB, instancia \texttt{g6.2xlarge}) en las fases iniciales de experimentación con un pico de uso de $\approx$20 GB; la NVIDIA L40S (48 GB, instancia \texttt{g6e.2xlarge}) en fases intermedias con un pico de $\approx$40 GB; y la RTX PRO 6000 (96 GB, instancia \texttt{g7e.2xlarge}) para el entrenamiento y benchmarking final, con un pico de 66,76 GB sobre los 300.000 samples con batch 64. Esta progresión garantiza que las métricas de latencia y throughput son representativas del entorno operacional real.

El checkpoint evaluado corresponde al modelo final almacenado en S3 (\texttt{s3://tfm-slm-benchmarks/benchmark\_hybrid.json}), descargado automáticamente por el sistema de benchmarking. El conjunto de validación, igualmente descargado desde S3, consiste en \textbf{300.000 muestras} del dataset mixto de diálogos e instrucciones, procesadas en lotes de 64 tokens, lo que resulta en 4.688 iteraciones de evaluación.

\begin{table}[H]
\centering
\begin{tabularx}{\textwidth}{|l|X|}
\hline
\textbf{Parámetro} & \textbf{Valor} \\ \hline
Modelo evaluado & HybridBlock (12 capas, 768 oculto, 12 cabezas) \\ \hline
Dispositivo & CUDA (NVIDIA RTX PRO 6000 Server, 96 GB VRAM) \\ \hline
Muestras de validación & 300.000 \\ \hline
Tamaño de lote & 64 \\ \hline
Iteraciones totales & 4.688 \\ \hline
Tiempo total de evaluación & 3.399,88 s ($\approx$56 min) \\ \hline
\end{tabularx}
\caption{Configuración del entorno de benchmarking.}
\label{tab:benchmark_config}
\end{table}

% ---------------------------------------------------------------
\subsection{Métricas de Evaluación}

El benchmark cubre diez métricas agrupadas en tres categorías: calidad del modelo, eficiencia computacional y huella de recursos.

\subsubsection{Calidad del Modelo}

\paragraph{Pérdida y Perplejidad.}
La \textbf{pérdida cross-entropy} final sobre el conjunto de validación es $\mathcal{L} = 1{,}091$. La \textbf{perplejidad}, definida como $\text{PPL} = e^{\mathcal{L}}$, alcanza un valor de \textbf{2,98}. La perplejidad puede interpretarse como el número efectivo de tokens igualmente probables que el modelo considera en cada posición: un valor de 2,98 indica que el modelo está altamente concentrado en sus predicciones, reduciendo la incertidumbre media a menos de 3 candidatos por token.

El log de evaluación muestra que la perplejidad se estabilizó en $\approx 2{,}98$ a partir del batch 70 (de 4.688), permaneciendo constante hasta el final. El siguiente extracto del fichero \texttt{.output/log\_benchmark.log} ilustra la progresión inicial y la rápida estabilización:

\begin{lstlisting}[language=bash, caption={Extracto de \texttt{log\_benchmark.log}: estabilización de perplejidad a partir del batch 70.}]
INFO: [10/4688]  Loss: 1.0668 | Perplexity: 2.87 | Acc@1: 80.5%
INFO: [20/4688]  Loss: 1.1011 | Perplexity: 2.91 | Acc@1: 80.3%
INFO: [30/4688]  Loss: 1.1393 | Perplexity: 2.92 | Acc@1: 80.2%
INFO: [50/4688]  Loss: 1.1912 | Perplexity: 2.96 | Acc@1: 79.9%
INFO: [70/4688]  Loss: 1.1335 | Perplexity: 2.98 | Acc@1: 79.8%
INFO: [80/4688]  Loss: 1.1482 | Perplexity: 2.98 | Acc@1: 79.8%
INFO: [100/4688] Loss: 1.1762 | Perplexity: 2.98 | Acc@1: 79.8%
INFO: [150/4688] Loss: 1.1127 | Perplexity: 2.98 | Acc@1: 79.7%
\end{lstlisting}

Esta estabilidad es indicativa de una distribución de validación uniforme y de un modelo bien generalizado sin sobreajuste localizado.

\paragraph{Precisión de Predicción.}
Se reportan tres niveles de precisión en la predicción del siguiente token:

\begin{table}[H]
\centering
\begin{tabularx}{\textwidth}{|l|c|X|}
\hline
\textbf{Métrica} & \textbf{Valor} & \textbf{Interpretación} \\ \hline
Top-1 Accuracy & 79,75\% & El token correcto es la primera predicción del modelo \\ \hline
Top-5 Accuracy & 91,14\% & El token correcto está entre las 5 primeras predicciones \\ \hline
Top-10 Accuracy & 93,56\% & El token correcto está entre las 10 primeras predicciones \\ \hline
\end{tabularx}
\caption{Métricas de precisión de predicción del modelo híbrido.}
\label{tab:accuracy}
\end{table}

Una precisión Top-1 del 79,75\% sobre 300.000 muestras es un resultado cuantitativamente sólido para un modelo entrenado desde cero con 167M de parámetros. La diferencia de 11,39 puntos porcentuales entre Top-1 y Top-5 refleja la naturaleza inherentemente ambigua del lenguaje natural: en muchos contextos, múltiples tokens son igualmente válidos como continuación, y el modelo los sitúa correctamente entre sus candidatos principales.

No obstante, estas métricas deben interpretarse con cautela. Las leyes de escala de Chinchilla~\cite{hoffmann2022chinchilla} establecen que un modelo de 167M de parámetros requiere aproximadamente \textbf{3.340 millones de tokens} para alcanzar capacidades generativas óptimas. Este proyecto procesó $\approx$\textbf{713 millones de tokens} de entrenamiento (46.439 muestras $\times$ 1.024 tokens $\times$ 15 épocas; las 300.000 muestras restantes del dataset se reservaron para esta evaluación) --- 4,7 veces menos del volumen óptimo (ratio 4,27:1 frente al 20:1 de Chinchilla). Una alta precisión de predicción sobre datos en distribución no garantiza coherencia en generación libre, ya que en inferencia autoregresiva los errores se propagan: un token mal predicho condiciona todos los siguientes. El análisis cualitativo completo se recoge en el Capítulo~\ref{sec:conclusiones}.

\subsubsection{Eficiencia Computacional}

\paragraph{Latencia e Inferencia.}
La latencia media por token en modo de inferencia individual es de \textbf{7,00 ms}, equivalente a una cadencia de \textbf{142,76 tokens/segundo} en modo secuencial. Este valor es determinante para aplicaciones conversacionales, donde se percibe como fluido cualquier modelo que supere los 50 tokens/segundo.

En modo de inferencia por lotes (\textit{batch inference}) se observan dos valores de throughput con naturalezas distintas. Durante la ejecución del benchmark, el log reporta de forma continua un rendimiento sostenido de $\approx$\textbf{146.000 tokens/segundo} (con un arranque inicial de 142.170 tok/s en el primer bloque). Sin embargo, el resumen final generado por el sistema de benchmarking registra un valor de \textbf{90.268 tokens/segundo} bajo la clave \texttt{batch\_throughput\_tokens\_per\_sec}.

Esta discrepancia se explica por la diferencia de granularidad en la medición: el log refleja el throughput instantáneo por ventana de 10 batches, mientras que el JSON reporta el throughput efectivo global, que incluye el tiempo de carga del checkpoint desde S3, la inicialización del dataset y la descarga de artefactos al inicio de la sesión. El valor operativo representativo del rendimiento de la GPU en inferencia pura es el del log ($\approx$146.000 tok/s), mientras que el del JSON refleja el coste real de un despliegue completo de principio a fin.

\subsubsection{Huella de Recursos}

\paragraph{Tamaño del Modelo.}
El modelo cuenta con \textbf{166.980.864 parámetros}, todos ellos entrenables, con un tamaño en disco de \textbf{636,98 MB}. Esto lo sitúa en la categoría de Small Language Model (SLM), equiparable en orden de magnitud a GPT-2 Medium (345M params) pero con una huella de memoria significativamente menor gracias a la arquitectura híbrida.

\paragraph{Memoria en Pico.}
El pico de memoria GPU durante la evaluación fue de \textbf{66,76 GB}, sobre los 96 GB disponibles en la NVIDIA RTX PRO 6000 Server. Este valor corresponde a 300.000 muestras con batch 64 en modo solo forward pass (sin gradientes), con activaciones intermedias acumuladas para el cómputo de métricas. Durante la inferencia en producción (un único token por paso con KV-cache, sin activaciones intermedias), el consumo desciende a $\approx$20 GB, compatible con la NVIDIA L4 empleada en fases anteriores del proyecto. El pico durante el entrenamiento completo se documenta en el Capítulo~\ref{sec:infraestructura}.

% ---------------------------------------------------------------
\subsection{Análisis de Estabilidad}

Una de las conclusiones más relevantes del log de benchmarking es la notable estabilidad de todas las métricas a lo largo de las 4.688 iteraciones. La perplejidad oscila en un rango de apenas $\pm 0{,}01$ tras la fase inicial de estabilización (primeros 200 batches), y la precisión Top-1 se mantiene entre 79,7\% y 79,8\% durante el 95\% de la evaluación.

Esta estabilidad no es trivial: indica que el modelo generaliza de forma homogénea a todos los subconjuntos del dataset de validación, sin degradación en muestras de dominios específicos ni varianza anómala que sugiera problemas de sobreajuste parcial.

Los picos de pérdida por batch observados en el log (valores ocasionales superiores a 1,2) son normales en evaluación de lenguaje natural y reflejan la presencia de muestras intrínsecamente difíciles (cambios abruptos de dominio, nombres propios infrecuentes, código), no un problema del modelo.

% ---------------------------------------------------------------
\subsection{Resumen de Resultados}

\begin{table}[H]
\centering
\begin{tabularx}{\textwidth}{|l|c|X|}
\hline
\textbf{Métrica} & \textbf{Resultado} & \textbf{Categoría} \\ \hline
Loss (cross-entropy) & 1,091 & Calidad \\ \hline
Perplejidad & 2,98 & Calidad \\ \hline
Top-1 Accuracy & 79,75\% & Calidad \\ \hline
Top-5 Accuracy & 91,14\% & Calidad \\ \hline
Top-10 Accuracy & 93,56\% & Calidad \\ \hline
Latencia por token & 7,00 ms & Eficiencia \\ \hline
Throughput (single token) & 142,76 tok/s & Eficiencia \\ \hline
Throughput (batch, GPU pura) & $\approx$146.000 tok/s & Eficiencia \\ \hline
Throughput (batch, despliegue completo) & 90.268 tok/s & Eficiencia \\ \hline
Parámetros totales & 166.980.864 & Recursos \\ \hline
Tamaño del modelo & 636,98 MB & Recursos \\ \hline
\end{tabularx}
\caption{Resumen de las 10 métricas del benchmark del modelo híbrido tfm-slm.}
\label{tab:benchmark_summary}
\end{table}

Los resultados cuantitativos son sólidos: una perplejidad de 2,98 y una precisión Top-1 del 79,75\% sobre 300.000 muestras son cifras competitivas para un modelo de 167M de parámetros entrenado desde cero. Sin embargo, estas métricas miden capacidad predictiva \textit{en distribución} ---sobre datos extraídos de la misma mezcla de entrenamiento---, no calidad de generación libre.

Como se documenta en el Capítulo~\ref{sec:conclusiones}, la evaluación cualitativa mediante inferencia real revela un cuadro distinto: el modelo produce bucles de repetición, incoherencia semántica y mezcla de idiomas de forma sistemática. La causa no es arquitectónica sino de escala de datos: con 713 millones de tokens de entrenamiento frente a los 3.340 millones que prescriben las leyes de Chinchilla~\cite{hoffmann2022chinchilla} para un modelo de este tamaño (ratio 4,27:1 frente al óptimo de 20:1, 4,7 veces menos), las capacidades generativas emergentes no se desarrollan. Las métricas del benchmark son necesarias pero no suficientes para evaluar la utilidad real de un modelo de lenguaje.